\documentclass[10pt,a4paper]{amsart}
\usepackage{amsmath,amssymb,amsthm,graphicx}
% Retained minimal preamble for the Stacks Project chapter varieties.tex.
% Verbatim excerpt of the chapter preamble: Xy-pic setup (native preamble lines 14-18), the native theorem
% declarations with the mathematical macros (native preamble lines 154-217), and the project own no-op
% definitions of the reference/slogan/history display environments (native wrapper.tex lines 1-3) that the
% counted statement uses for its slogan. Copied mechanically; nothing added, removed or redefined.
% For commutative diagrams we use Xy-pic
\usepackage[all]{xy}

% We use 2cell for 2-commutative diagrams.
\xyoption{2cell}
% Theorem environments.
%
\theoremstyle{plain}
\newtheorem{theorem}[subsection]{Theorem}
\newtheorem{proposition}[subsection]{Proposition}
\newtheorem{lemma}[subsection]{Lemma}

\theoremstyle{definition}
\newtheorem{definition}[subsection]{Definition}
\newtheorem{example}[subsection]{Example}
\newtheorem{exercise}[subsection]{Exercise}
\newtheorem{situation}[subsection]{Situation}

\theoremstyle{remark}
\newtheorem{remark}[subsection]{Remark}
\newtheorem{remarks}[subsection]{Remarks}

\numberwithin{equation}{subsection}

% Macros
%
\def\lim{\mathop{\mathrm{lim}}\nolimits}
\def\colim{\mathop{\mathrm{colim}}\nolimits}
\def\Spec{\mathop{\mathrm{Spec}}}
\def\Hom{\mathop{\mathrm{Hom}}\nolimits}
\def\Ext{\mathop{\mathrm{Ext}}\nolimits}
\def\SheafHom{\mathop{\mathcal{H}\!\mathit{om}}\nolimits}
\def\SheafExt{\mathop{\mathcal{E}\!\mathit{xt}}\nolimits}
\def\Sch{\mathit{Sch}}
\def\Mor{\mathop{\mathrm{Mor}}\nolimits}
\def\Ob{\mathop{\mathrm{Ob}}\nolimits}
\def\Sh{\mathop{\mathit{Sh}}\nolimits}
\def\NL{\mathop{N\!L}\nolimits}
\def\CH{\mathop{\mathrm{CH}}\nolimits}
\def\proetale{{pro\text{-}\acute{e}tale}}
\def\etale{{\acute{e}tale}}
\def\QCoh{\mathit{QCoh}}
\def\Ker{\mathop{\mathrm{Ker}}}
\def\Im{\mathop{\mathrm{Im}}}
\def\Coker{\mathop{\mathrm{Coker}}}
\def\Coim{\mathop{\mathrm{Coim}}}

% Boxtimes
%
\DeclareMathSymbol{\boxtimes}{\mathbin}{AMSa}{"02}

%
% Macros for moduli stacks/spaces
%
\def\QCohstack{\mathcal{QC}\!\mathit{oh}}
\def\Cohstack{\mathcal{C}\!\mathit{oh}}
\def\Spacesstack{\mathcal{S}\!\mathit{paces}}
\def\Quotfunctor{\mathrm{Quot}}
\def\Hilbfunctor{\mathrm{Hilb}}
\def\Curvesstack{\mathcal{C}\!\mathit{urves}}
\def\Polarizedstack{\mathcal{P}\!\mathit{olarized}}
\def\Complexesstack{\mathcal{C}\!\mathit{omplexes}}
% \Pic is the operator that assigns to X its picard group, usage \Pic(X)
% \Picardstack_{X/B} denotes the Picard stack of X over B
% \Picardfunctor_{X/B} denotes the Picard functor of X over B
\def\Pic{\mathop{\mathrm{Pic}}\nolimits}
\def\Picardstack{\mathcal{P}\!\mathit{ic}}
\def\Picardfunctor{\mathrm{Pic}}
\def\Deformationcategory{\mathcal{D}\!\mathit{ef}}
\newenvironment{reference}{}{}
\newenvironment{slogan}{}{}
\newenvironment{history}{}{}
\usepackage{cleveref}
\crefname{originallemma}{original lemma}{original lemmas}
\Crefname{originallemma}{Original lemma}{Original lemmas}
\makeatletter
\newlabel{fields-lemma-common-extension-field}{{0H7K}{source}{}{sourcepointer-1}{}}
\newlabel{fields-lemma-common-extension-field@cref}{{[originallemma][0][]0H7K}{[1][0][]source}{}{}{}}
\makeatother
\crefname{originallemma}{original lemma}{original lemmas}
\Crefname{originallemma}{Original lemma}{Original lemmas}
\makeatletter
\newlabel{algebra-lemma-descent-reduced}{{033F}{source}{}{sourcepointer-2}{}}
\newlabel{algebra-lemma-descent-reduced@cref}{{[originallemma][0][]033F}{[1][0][]source}{}{}{}}
\makeatother
\crefname{originallemma}{original lemma}{original lemmas}
\Crefname{originallemma}{Original lemma}{Original lemmas}
\makeatletter
\newlabel{lemma-check-only-finite-inseparable-extensions}{{035Y}{source}{}{sourcepointer-3}{}}
\newlabel{lemma-check-only-finite-inseparable-extensions@cref}{{[originallemma][0][]035Y}{[1][0][]source}{}{}{}}
\makeatother
\title{Geometric reducedness at a point may be checked after an arbitrary field extension}
\author{The Stacks Project Authors}
\date{stacks-project:varieties.tex:local-snapshot-20261011; GNU FDL 1.2 or later}
\begin{document}
\maketitle
\noindent\textit{Editorial note.} All delivered mathematical text is the author's own, copied from the pinned Stacks Project chapter file varieties.tex; no statement or proof step was rewritten, reordered or removed.\par\medskip
\noindent\textit{Editorial note.} The preamble is the chapter's own retained minimal preamble (sources/stacks-varieties/retained-preamble.tex), a verbatim excerpt of the native Xy-pic setup, of the native theorem declarations with mathematical macros, and of the project's own no-op definitions of the reference display environments; no macro was invented or redefined.\par\medskip
\noindent\textit{Editorial note.} The same-chapter reference is located at its \textbackslash{}label line inside the pinned chapter file; the two cross-chapter pointers are located in the project's own tag list staged in the snapshot (sources/stacks-chow/tags.txt), at the verbatim lines carrying their canonical tags.\par\medskip
\noindent\textit{Allowed standard prerequisites (editorial).}\begin{itemize}
\item Geometric reducedness at a point, tested after finite purely inseparable extensions. Same chapter, Lemma labelled lemma-check-only-finite-inseparable-extensions (source line 463, Stacks tag 035Y); invoked at proof lines 520-521. Supplies that checking reducedness of the local ring after every finite purely inseparable extension is enough; the transport of that condition from k'' back to x is the counted content.
\item Descent of reducedness along a faithfully flat local ring map. Algebra, Lemma labelled algebra-lemma-descent-reduced (canonical Stacks tag 033F); invoked at proof lines 518-519. Supplies that reducedness of the bigger local ring implies reducedness of the smaller one.
\item Existence of a common extension of two field extensions. Fields, Lemma labelled fields-lemma-common-extension-field (canonical Stacks tag 0H7K); invoked at proof lines 509-510. Supplies the field k''' dominating both k'' and k', over which the two base changes may be compared.
\end{itemize}\medskip
\ifdefined\hypertarget\hypertarget{sourcepointer-1}{}\fi
\noindent\textit{Source reference.} Original lemma 0H7K: sources/stacks-chow/tags.txt, source lines 21460--21460. Cross-chapter standard prerequisite invoked once at proof lines 509-510 for the existence of the common extension k'''; the faithful flatness and descent arguments are retained in the counted span.\par\medskip
\ifdefined\hypertarget\hypertarget{sourcepointer-2}{}\fi
\noindent\textit{Source reference.} Original lemma 033F: sources/stacks-chow/tags.txt, source lines 4054--4054. Cross-chapter standard prerequisite invoked once at proof lines 518-519 for descent of reducedness along a faithfully flat local ring map; the identification of the ring map as faithfully flat is retained in the counted span.\par\medskip
\ifdefined\hypertarget\hypertarget{sourcepointer-3}{}\fi
\noindent\textit{Source reference.} Original lemma 035Y: sources/stacks-varieties/pinned-original-source.tex, source lines 463--463. Same-chapter pointer invoked at proof lines 520-521 for the criterion that finite purely inseparable test extensions suffice; the transport of reducedness to each such extension is retained in the counted span.\par\medskip
\begin{lemma}
\label{lemma-geometrically-reduced-upstairs}
Let $k$ be a field.
Let $X$ be a scheme over $k$.
Let $k'/k$ be a field extension.
Let $x \in X$ be a point, and let $x' \in X_{k'}$ be a point lying over $x$.
The following are equivalent
\begin{enumerate}
\item $X$ is geometrically reduced at $x$,
\item $X_{k'}$ is geometrically reduced at $x'$.
\end{enumerate}
In particular, $X$ is geometrically reduced over $k$ if and only if
$X_{k'}$ is geometrically reduced over $k'$.
\end{lemma}

\begin{proof}
It is clear that (1) implies (2). Assume (2).
Let $k''/k$ be a finite purely inseparable field extension
and let $x'' \in X_{k''}$ be a point lying over $x$ (actually it is
unique). Choose a common field extension $k'''/k$ of $k''/k$ and $k'/k$, see
Fields, Lemma \ref{fields-lemma-common-extension-field}.
Choose a point $x''' \in X_{k'''}$ lying over both $x'$ and $x''$.
Consider the map of local rings
$$
\mathcal{O}_{X_{k''}, x''} \longrightarrow \mathcal{O}_{X_{k'''}, x'''}.
$$
This is a flat local ring homomorphism and hence faithfully flat.
By (2) we see that the local ring on the right is reduced.
Thus by Algebra, Lemma \ref{algebra-lemma-descent-reduced}
we conclude that $\mathcal{O}_{X_{k''}, x''}$ is reduced.
Thus by Lemma \ref{lemma-check-only-finite-inseparable-extensions}
we conclude that $X$ is geometrically reduced at $x$.
\end{proof}

\end{document}
